\section{}

\paragraph{1114年}\eventanchor{JIN-1114-EVENT-078}{JIN-1114-EVENT-078}　完颜阿骨打起兵反辽。此项以《金史》〈本纪〉、〈交聘表〉及相关志；《宋史》与《建炎以来系年要录》相关条 为基本来源线索；编年与官修材料能够确认年序和制度用语，但对地方执行、人数、动机或社会后果的判断仍须同文书、碑刻、考古及对方叙事互校。

\paragraph{1115年}\eventanchor{JIN-1115-EVENT-079}{JIN-1115-EVENT-079}　完颜阿骨打称帝建金。此项以《金史》〈本纪〉、〈交聘表〉及相关志；《宋史》与《建炎以来系年要录》相关条 为基本来源线索；编年与官修材料能够确认年序和制度用语，但对地方执行、人数、动机或社会后果的判断仍须同文书、碑刻、考古及对方叙事互校。

\paragraph{1115年}\eventanchor{JIN-1115-REGNAL-YEAR}{JIN-1115-REGNAL-YEAR}　金以金太祖在位第1年作为诏令、赋役与外交文书的纪年框架。此项以《金史》〈本纪〉、〈交聘表〉及相关志；《宋史》与《建炎以来系年要录》相关条 为基本来源线索；编年与官修材料能够确认年序和制度用语，但对地方执行、人数、动机或社会后果的判断仍须同文书、碑刻、考古及对方叙事互校。

\paragraph{1116年}\eventanchor{JIN-1116-REGNAL-YEAR}{JIN-1116-REGNAL-YEAR}　金以金太祖在位第2年作为诏令、赋役与外交文书的纪年框架。此项以《金史》〈本纪〉、〈交聘表〉及相关志；《宋史》与《建炎以来系年要录》相关条 为基本来源线索；编年与官修材料能够确认年序和制度用语，但对地方执行、人数、动机或社会后果的判断仍须同文书、碑刻、考古及对方叙事互校。

\paragraph{1117年}\eventanchor{JIN-1117-REGNAL-YEAR}{JIN-1117-REGNAL-YEAR}　金以金太祖在位第3年作为诏令、赋役与外交文书的纪年框架。此项以《金史》〈本纪〉、〈交聘表〉及相关志；《宋史》与《建炎以来系年要录》相关条 为基本来源线索；编年与官修材料能够确认年序和制度用语，但对地方执行、人数、动机或社会后果的判断仍须同文书、碑刻、考古及对方叙事互校。

\paragraph{1118年}\eventanchor{JIN-1118-EVENT-080}{JIN-1118-EVENT-080}　金与宋讨论共同攻辽。此项以《金史》〈本纪〉、〈交聘表〉及相关志；《宋史》与《建炎以来系年要录》相关条 为基本来源线索；编年与官修材料能够确认年序和制度用语，但对地方执行、人数、动机或社会后果的判断仍须同文书、碑刻、考古及对方叙事互校。

\paragraph{1118年}\eventanchor{JIN-1118-REGNAL-YEAR}{JIN-1118-REGNAL-YEAR}　金以金太祖在位第4年作为诏令、赋役与外交文书的纪年框架。此项以《金史》〈本纪〉、〈交聘表〉及相关志；《宋史》与《建炎以来系年要录》相关条 为基本来源线索；编年与官修材料能够确认年序和制度用语，但对地方执行、人数、动机或社会后果的判断仍须同文书、碑刻、考古及对方叙事互校。

\paragraph{1119年}\eventanchor{JIN-1119-REGNAL-YEAR}{JIN-1119-REGNAL-YEAR}　金以金太祖在位第5年作为诏令、赋役与外交文书的纪年框架。此项以《金史》〈本纪〉、〈交聘表〉及相关志；《宋史》与《建炎以来系年要录》相关条 为基本来源线索；编年与官修材料能够确认年序和制度用语，但对地方执行、人数、动机或社会后果的判断仍须同文书、碑刻、考古及对方叙事互校。

\paragraph{1120年}\eventanchor{JIN-1120-EVENT-081}{JIN-1120-EVENT-081}　金宋达成海上之盟。此项以《金史》〈本纪〉、〈交聘表〉及相关志；《宋史》与《建炎以来系年要录》相关条 为基本来源线索；编年与官修材料能够确认年序和制度用语，但对地方执行、人数、动机或社会后果的判断仍须同文书、碑刻、考古及对方叙事互校。

\paragraph{1120年}\eventanchor{JIN-1120-REGNAL-YEAR}{JIN-1120-REGNAL-YEAR}　金以金太祖在位第6年作为诏令、赋役与外交文书的纪年框架。此项以《金史》〈本纪〉、〈交聘表〉及相关志；《宋史》与《建炎以来系年要录》相关条 为基本来源线索；编年与官修材料能够确认年序和制度用语，但对地方执行、人数、动机或社会后果的判断仍须同文书、碑刻、考古及对方叙事互校。

\paragraph{1121年}\eventanchor{JIN-1121-REGNAL-YEAR}{JIN-1121-REGNAL-YEAR}　金以金太祖在位第7年作为诏令、赋役与外交文书的纪年框架。此项以《金史》〈本纪〉、〈交聘表〉及相关志；《宋史》与《建炎以来系年要录》相关条 为基本来源线索；编年与官修材料能够确认年序和制度用语，但对地方执行、人数、动机或社会后果的判断仍须同文书、碑刻、考古及对方叙事互校。

\paragraph{1122年}\eventanchor{JIN-1122-EVENT-082}{JIN-1122-EVENT-082}　金攻取辽南京燕京。此项以《金史》〈本纪〉、〈交聘表〉及相关志；《宋史》与《建炎以来系年要录》相关条 为基本来源线索；编年与官修材料能够确认年序和制度用语，但对地方执行、人数、动机或社会后果的判断仍须同文书、碑刻、考古及对方叙事互校。

\paragraph{1122年}\eventanchor{JIN-1122-REGNAL-YEAR}{JIN-1122-REGNAL-YEAR}　金以金太祖在位第8年作为诏令、赋役与外交文书的纪年框架。此项以《金史》〈本纪〉、〈交聘表〉及相关志；《宋史》与《建炎以来系年要录》相关条 为基本来源线索；编年与官修材料能够确认年序和制度用语，但对地方执行、人数、动机或社会后果的判断仍须同文书、碑刻、考古及对方叙事互校。

\paragraph{1123年}\eventanchor{JIN-1123-EVENT-083}{JIN-1123-EVENT-083}　金将燕京部分地区交予宋并要求岁币。此项以《金史》〈本纪〉、〈交聘表〉及相关志；《宋史》与《建炎以来系年要录》相关条 为基本来源线索；编年与官修材料能够确认年序和制度用语，但对地方执行、人数、动机或社会后果的判断仍须同文书、碑刻、考古及对方叙事互校。

\paragraph{1123年}\eventanchor{JIN-1123-REGNAL-YEAR}{JIN-1123-REGNAL-YEAR}　金以金太祖在位第9年作为诏令、赋役与外交文书的纪年框架。此项以《金史》〈本纪〉、〈交聘表〉及相关志；《宋史》与《建炎以来系年要录》相关条 为基本来源线索；编年与官修材料能够确认年序和制度用语，但对地方执行、人数、动机或社会后果的判断仍须同文书、碑刻、考古及对方叙事互校。

\paragraph{1124年}\eventanchor{JIN-1124-REGNAL-YEAR}{JIN-1124-REGNAL-YEAR}　金以金太宗在位第1年作为诏令、赋役与外交文书的纪年框架。此项以《金史》〈本纪〉、〈交聘表〉及相关志；《宋史》与《建炎以来系年要录》相关条 为基本来源线索；编年与官修材料能够确认年序和制度用语，但对地方执行、人数、动机或社会后果的判断仍须同文书、碑刻、考古及对方叙事互校。

\paragraph{1125年}\eventanchor{JIN-1125-EVENT-084}{JIN-1125-EVENT-084}　金军俘获辽天祚帝，辽朝覆亡。此项以《金史》〈本纪〉、〈交聘表〉及相关志；《宋史》与《建炎以来系年要录》相关条 为基本来源线索；编年与官修材料能够确认年序和制度用语，但对地方执行、人数、动机或社会后果的判断仍须同文书、碑刻、考古及对方叙事互校。

\paragraph{1125年}\eventanchor{JIN-1125-EVENT-085}{JIN-1125-EVENT-085}　金军分路南侵宋境。此项以《金史》〈本纪〉、〈交聘表〉及相关志；《宋史》与《建炎以来系年要录》相关条 为基本来源线索；编年与官修材料能够确认年序和制度用语，但对地方执行、人数、动机或社会后果的判断仍须同文书、碑刻、考古及对方叙事互校。

\paragraph{1125年}\eventanchor{JIN-1125-REGNAL-YEAR}{JIN-1125-REGNAL-YEAR}　金以金太宗在位第2年作为诏令、赋役与外交文书的纪年框架。此项以《金史》〈本纪〉、〈交聘表〉及相关志；《宋史》与《建炎以来系年要录》相关条 为基本来源线索；编年与官修材料能够确认年序和制度用语，但对地方执行、人数、动机或社会后果的判断仍须同文书、碑刻、考古及对方叙事互校。

\paragraph{1126年}\eventanchor{JIN-1126-EVENT-086}{JIN-1126-EVENT-086}　金军第一次围汴京并迫使宋廷议和。此项以《金史》〈本纪〉、〈交聘表〉及相关志；《宋史》与《建炎以来系年要录》相关条 为基本来源线索；编年与官修材料能够确认年序和制度用语，但对地方执行、人数、动机或社会后果的判断仍须同文书、碑刻、考古及对方叙事互校。

\paragraph{1126年}\eventanchor{JIN-1126-REGNAL-YEAR}{JIN-1126-REGNAL-YEAR}　金以金太宗在位第3年作为诏令、赋役与外交文书的纪年框架。此项以《金史》〈本纪〉、〈交聘表〉及相关志；《宋史》与《建炎以来系年要录》相关条 为基本来源线索；编年与官修材料能够确认年序和制度用语，但对地方执行、人数、动机或社会后果的判断仍须同文书、碑刻、考古及对方叙事互校。

\paragraph{1127年}\eventanchor{JIN-1127-EVENT-087}{JIN-1127-EVENT-087}　金军攻陷汴京，俘徽、钦二帝。此项以《金史》〈本纪〉、〈交聘表〉及相关志；《宋史》与《建炎以来系年要录》相关条 为基本来源线索；编年与官修材料能够确认年序和制度用语，但对地方执行、人数、动机或社会后果的判断仍须同文书、碑刻、考古及对方叙事互校。

\paragraph{1127年}\eventanchor{JIN-1127-REGNAL-YEAR}{JIN-1127-REGNAL-YEAR}　金以金太宗在位第4年作为诏令、赋役与外交文书的纪年框架。此项以《金史》〈本纪〉、〈交聘表〉及相关志；《宋史》与《建炎以来系年要录》相关条 为基本来源线索；编年与官修材料能够确认年序和制度用语，但对地方执行、人数、动机或社会后果的判断仍须同文书、碑刻、考古及对方叙事互校。

\paragraph{1128年}\eventanchor{JIN-1128-REGNAL-YEAR}{JIN-1128-REGNAL-YEAR}　金以金太宗在位第5年作为诏令、赋役与外交文书的纪年框架。此项以《金史》〈本纪〉、〈交聘表〉及相关志；《宋史》与《建炎以来系年要录》相关条 为基本来源线索；编年与官修材料能够确认年序和制度用语，但对地方执行、人数、动机或社会后果的判断仍须同文书、碑刻、考古及对方叙事互校。

\paragraph{1129年}\eventanchor{JIN-1129-EVENT-088}{JIN-1129-EVENT-088}　完颜宗弼等军深入江淮。此项以《金史》〈本纪〉、〈交聘表〉及相关志；《宋史》与《建炎以来系年要录》相关条 为基本来源线索；编年与官修材料能够确认年序和制度用语，但对地方执行、人数、动机或社会后果的判断仍须同文书、碑刻、考古及对方叙事互校。

\paragraph{1129年}\eventanchor{JIN-1129-REGNAL-YEAR}{JIN-1129-REGNAL-YEAR}　金以金太宗在位第6年作为诏令、赋役与外交文书的纪年框架。此项以《金史》〈本纪〉、〈交聘表〉及相关志；《宋史》与《建炎以来系年要录》相关条 为基本来源线索；编年与官修材料能够确认年序和制度用语，但对地方执行、人数、动机或社会后果的判断仍须同文书、碑刻、考古及对方叙事互校。

\paragraph{1130年}\eventanchor{JIN-1130-REGNAL-YEAR}{JIN-1130-REGNAL-YEAR}　金以金太宗在位第7年作为诏令、赋役与外交文书的纪年框架。此项以《金史》〈本纪〉、〈交聘表〉及相关志；《宋史》与《建炎以来系年要录》相关条 为基本来源线索；编年与官修材料能够确认年序和制度用语，但对地方执行、人数、动机或社会后果的判断仍须同文书、碑刻、考古及对方叙事互校。

\paragraph{1131年}\eventanchor{JIN-1131-EVENT-089}{JIN-1131-EVENT-089}　金扶植刘豫政权经营中原。此项以《金史》〈本纪〉、〈交聘表〉及相关志；《宋史》与《建炎以来系年要录》相关条 为基本来源线索；编年与官修材料能够确认年序和制度用语，但对地方执行、人数、动机或社会后果的判断仍须同文书、碑刻、考古及对方叙事互校。

\paragraph{1131年}\eventanchor{JIN-1131-REGNAL-YEAR}{JIN-1131-REGNAL-YEAR}　金以金太宗在位第8年作为诏令、赋役与外交文书的纪年框架。此项以《金史》〈本纪〉、〈交聘表〉及相关志；《宋史》与《建炎以来系年要录》相关条 为基本来源线索；编年与官修材料能够确认年序和制度用语，但对地方执行、人数、动机或社会后果的判断仍须同文书、碑刻、考古及对方叙事互校。

\paragraph{1132年}\eventanchor{JIN-1132-REGNAL-YEAR}{JIN-1132-REGNAL-YEAR}　金以金太宗在位第9年作为诏令、赋役与外交文书的纪年框架。此项以《金史》〈本纪〉、〈交聘表〉及相关志；《宋史》与《建炎以来系年要录》相关条 为基本来源线索；编年与官修材料能够确认年序和制度用语，但对地方执行、人数、动机或社会后果的判断仍须同文书、碑刻、考古及对方叙事互校。

\paragraph{1133年}\eventanchor{JIN-1133-REGNAL-YEAR}{JIN-1133-REGNAL-YEAR}　金以金太宗在位第10年作为诏令、赋役与外交文书的纪年框架。此项以《金史》〈本纪〉、〈交聘表〉及相关志；《宋史》与《建炎以来系年要录》相关条 为基本来源线索；编年与官修材料能够确认年序和制度用语，但对地方执行、人数、动机或社会后果的判断仍须同文书、碑刻、考古及对方叙事互校。

\paragraph{1134年}\eventanchor{JIN-1134-REGNAL-YEAR}{JIN-1134-REGNAL-YEAR}　金以金太宗在位第11年作为诏令、赋役与外交文书的纪年框架。此项以《金史》〈本纪〉、〈交聘表〉及相关志；《宋史》与《建炎以来系年要录》相关条 为基本来源线索；编年与官修材料能够确认年序和制度用语，但对地方执行、人数、动机或社会后果的判断仍须同文书、碑刻、考古及对方叙事互校。

\paragraph{1135年}\eventanchor{JIN-1135-REGNAL-YEAR}{JIN-1135-REGNAL-YEAR}　金以金太宗在位第12年作为诏令、赋役与外交文书的纪年框架。此项以《金史》〈本纪〉、〈交聘表〉及相关志；《宋史》与《建炎以来系年要录》相关条 为基本来源线索；编年与官修材料能够确认年序和制度用语，但对地方执行、人数、动机或社会后果的判断仍须同文书、碑刻、考古及对方叙事互校。

\paragraph{1136年}\eventanchor{JIN-1136-REGNAL-YEAR}{JIN-1136-REGNAL-YEAR}　金以金熙宗在位第1年作为诏令、赋役与外交文书的纪年框架。此项以《金史》〈本纪〉、〈交聘表〉及相关志；《宋史》与《建炎以来系年要录》相关条 为基本来源线索；编年与官修材料能够确认年序和制度用语，但对地方执行、人数、动机或社会后果的判断仍须同文书、碑刻、考古及对方叙事互校。

\paragraph{1137年}\eventanchor{JIN-1137-EVENT-090}{JIN-1137-EVENT-090}　金废除刘豫政权，改为直接经营中原。此项以《金史》〈本纪〉、〈交聘表〉及相关志；《宋史》与《建炎以来系年要录》相关条 为基本来源线索；编年与官修材料能够确认年序和制度用语，但对地方执行、人数、动机或社会后果的判断仍须同文书、碑刻、考古及对方叙事互校。

